\documentclass{amsart}
\usepackage{amsmath,amssymb,amsthm}

\title{CRT Decomposition of $\mathbb{Z}/10\mathbb{Z}$}
\author{Brayden Ross Sanders / 7Site LLC}
\date{2026}

\newtheorem{theorem}{Theorem}
\newtheorem{corollary}[theorem]{Corollary}
\newtheorem{definition}[theorem]{Definition}

\begin{document}
\maketitle

\begin{abstract}
We establish the canonical Chinese Remainder Theorem decomposition
$\mathbb{Z}/10\mathbb{Z} \cong \mathbb{Z}/2\mathbb{Z} \times \mathbb{Z}/5\mathbb{Z}$
and derive the unique Fourier embedding $v: \mathbb{Z}/10\mathbb{Z} \to \mathbb{R}^5$
that factors through this isomorphism.
\end{abstract}

\section{The CRT Isomorphism}

\begin{theorem}[CRT Isomorphism]
There is a canonical ring isomorphism
\[
  \varphi: \mathbb{Z}/10\mathbb{Z} \xrightarrow{\;\sim\;} \mathbb{Z}/2\mathbb{Z} \times \mathbb{Z}/5\mathbb{Z},
  \quad \varphi(\mathrm{op}) = (\mathrm{op} \bmod 2,\; \mathrm{op} \bmod 5).
\]
\end{theorem}

\begin{proof}
Since $\gcd(2,5)=1$ and $2 \cdot 5 = 10$, the CRT gives a ring isomorphism. The map
$\varphi$ sends each residue class to its pair of residues mod 2 and mod 5.
\end{proof}

\section{The Fourier Embedding}

\begin{definition}[CRT Fourier Embedding]
Write $\varphi(\mathrm{op}) = (\varepsilon, y)$ with $\varepsilon \in \{0,1\}$ and $y \in \{0,1,2,3,4\}$.
The \emph{CRT Fourier embedding} is the map $v: \mathbb{Z}/10\mathbb{Z} \to \mathbb{R}^5$:
\[
  v(\mathrm{op}) = \bigl(\varepsilon,\; \cos(2\pi y/5),\; \sin(2\pi y/5),\;
                        \cos(4\pi y/5),\; \sin(4\pi y/5)\bigr).
\]
\end{definition}

\begin{theorem}[Uniqueness of the Fourier Embedding]
$v$ is the unique embedding that (i) factors through the CRT isomorphism,
(ii) uses the standard Fourier basis of $\mathbb{Z}/5\mathbb{Z}$, and
(iii) maps the $\mathbb{Z}/2\mathbb{Z}$ flag to a single real coordinate.
\end{theorem}

\begin{proof}
The CRT isomorphism is unique. The standard Fourier basis of $\mathbb{Z}/5\mathbb{Z}$
consists of $\{\cos(2\pi k y/5), \sin(2\pi k y/5)\}$ for $k=1,2$, giving four real
components. Any other basis is a rotation of this one; the Fourier basis is canonical.
The $\mathbb{Z}/2\mathbb{Z}$ flag maps to $\{0,1\} \subset \mathbb{R}$ uniquely.
This yields $1+4=5$ dimensions total.
\end{proof}

\begin{corollary}[5D Force Vector is Derived]
The five dimensions (aperture $\leftrightarrow \varepsilon$,
pressure $\leftrightarrow \cos(2\pi y/5)$,
depth $\leftrightarrow \sin(2\pi y/5)$,
binding $\leftrightarrow \cos(4\pi y/5)$,
continuity $\leftrightarrow \sin(4\pi y/5)$)
are algebraic consequences of the ring structure of $\mathbb{Z}/10\mathbb{Z}$,
not empirical assignments.
\end{corollary}

\end{document}
